\documentclass[11pt,a4paper]{article}
\usepackage[margin=0.55in, top=0.5in, bottom=0.5in]{geometry}
\usepackage[utf8]{inputenc}
\usepackage[T1]{fontenc}
\usepackage{charter}
\usepackage{fontawesome5}
\usepackage{xcolor}
\usepackage{titlesec}
\usepackage{enumitem}
\usepackage{tikz}
\usepackage{hyperref}
\usepackage{tcolorbox}
\tcbuselibrary{skins, breakable}
\usepackage{microtype}

\definecolor{headerbg}{HTML}{FFFFFF}
\definecolor{accent}{HTML}{1E3A5F}     % Professional Dark Navy
\definecolor{accentblue}{HTML}{2B6CB0} % Slate Blue
\definecolor{lightgray}{HTML}{F3F4F6}
\definecolor{darktext}{HTML}{1F2937}
\definecolor{skillbg}{HTML}{374151}
\definecolor{bulletline}{HTML}{9CA3AF}
\definecolor{subtletext}{HTML}{4B5563}

\hypersetup{colorlinks=true, linkcolor=accentblue, urlcolor=accentblue}

\titleformat{\section}{\large\bfseries\color{accent}}{}{0em}{}[\vspace{1pt}\titlerule]
\titlespacing{\section}{0pt}{11pt}{7pt}

\newtcbox{\skillbadge}{
    on line, box align=base,
    colback=skillbg, colframe=skillbg,
    colupper=white,
    fontupper=\small\bfseries,
    arc=4pt, boxrule=0pt,
    left=6pt, right=6pt, top=3pt, bottom=3pt,
    boxsep=0pt
}

% Caja de entrada de experiencia con borde izquierdo elegante y continuo
\newtcolorbox{timelineentry}{
    blanker,
    breakable,
    borderline west={1.5pt}{0pt}{bulletline},
    left=1.2em,
    top=0pt,
    bottom=0pt,
    before skip=0.8ex,
    after skip=2.6ex
}

\newcommand{\cventry}[5]{%
  \noindent
  \begin{minipage}{\textwidth}
    \vspace{0pt}%
    \tikz[baseline=-0.7ex]{\node[draw=accent, fill=white, circle, inner sep=2.5pt, line width=1.5pt] {};}%
    \hspace{0.5em}%
    {\large\bfseries\color{accent} #1}%
    \hfill{\small\bfseries\color{accentblue} #4}\par
    \vspace{0.2ex}%
    \hspace{1.5em}{\normalsize\bfseries\color{darktext} #2}%
    \hfill{\small\color{subtletext}\textit{#3}}\par
  \end{minipage}\par
  \vspace{0.4ex}%
  \begin{timelineentry}
    \small\color{darktext}
    \begin{itemize}[label={\color{accentblue}\raisebox{0.2ex}{\tiny\faCircle}}, leftmargin=1.2em, itemsep=0.4ex, parsep=0pt, topsep=0pt, partopsep=0pt]
      #5
    \end{itemize}
  \end{timelineentry}
}

\begin{document}
\thispagestyle{empty}
\color{darktext}

% Encabezado
\begin{center}
    {\Huge \bfseries\color{accent} SANTIAGO ALEJANDRO OCHOA QUESADA} \\[6pt]
    {\small\color{darktext}
    \faEnvelope\ \href{mailto:santy.alejo4@gmail.com}{santy.alejo4@gmail.com}
    \quad\textcolor{bulletline}{\textbar}\quad
    \faPhone\ +57 320 351 7153
    \quad\textcolor{bulletline}{\textbar}\quad
    \faMapMarker*\ Bogotá, Colombia \\[4pt]
    \faLinkedin\ \href{https://linkedin.com/in/santiago-alejandro-ochoa-quesada-a07a18288/}{LinkedIn}
    \quad\textcolor{bulletline}{\textbar}\quad
    \faGithub\ \href{https://github.com/Sochoaq}{GitHub}
    }
\end{center}
\vspace{4pt}

\section{Perfil Profesional}
\textbf{Ingeniero Electrónico} de la Universidad Nacional de Colombia (promedio 4.6/5.0) y actual \textbf{Ingeniero de Proyectos} en AM4. Especializado en AI Engineering, desarrollo backend en Python y construcción de pipelines de datos y automatización inteligente. Enfocado en el diseño e implementación de sistemas de IA Generativa, arquitecturas RAG, agentes autónomos y pipelines de MLOps con monitoreo continuo. Experiencia integrando APIs REST, entornos analíticos cloud/Databricks y orquestación con Docker. Nivel de inglés B2 profesional para comunicación y gestión técnica en equipos multidisciplinarios.

\section{Experiencia Profesional}

\cventry{Ingeniero de Proyectos (Desarrollo e IA)}{AM4}{Bogotá}{04/2026 - Actualidad}
{
    \item Diseño y desarrollo de soluciones de automatización inteligente y pipelines en Python para el procesamiento, limpieza y análisis de datos de telemetría y variables complejas.
    \item Implementación de flujos de analítica de datos e ingestión estructurada en entornos analíticos y Databricks, optimizando la preparación y consulta de grandes volúmenes de datos operativos.
    \item Construcción de microservicios y APIs REST para integrar modelos analíticos y de visión artificial con plataformas cliente-servidor y bases de datos SQL.
    \item Implementación de buenas prácticas de MLOps y versionado de datos, estableciendo métricas de evaluación, trazabilidad y monitoreo continuo de señales.
    \item Despliegue de herramientas contenerizadas en Docker y aseguramiento de calidad con pruebas automatizadas (Pytest), alcanzando 98\% de confiabilidad en campo para clientes industriales como Refoenergy, Brinsa y Geopark.
    \item Diagnóstico e instrumentación de sistemas críticos, reduciendo tiempos de validación técnica y estructurando la captura consistente de datos.
    \item Coordinación técnica de entregables, elaboración de reportes ejecutivos y comunicación efectiva de avances e hitos con stakeholders y directores de proyecto.
}

\cventry{Auxiliar Estudiantil}{Universidad Nacional de Colombia}{Bogotá}{07/2024 - 11/2025}
{
    \item Apoyo en laboratorios mediante organización, validación y análisis de resultados experimentales, programación, simulación y depuración de modelos numéricos y señales.
    \item Preparación y verificación de prácticas, registros, equipos de instrumentación y material técnico, asegurando consistencia en la toma de datos durante las sesiones.
    \item Acompañamiento a estudiantes en análisis de datos, identificación de errores y resolución de problemas de programación y sistemas digitales.
    \item Elaboración de guías, reportes y presentaciones técnicas para documentar procedimientos, fallas, hallazgos y aprendizajes de manera estructurada.
}

\section{Educación}

\cventry{Ingeniería Electrónica}{Universidad Nacional de Colombia}{Bogotá}{02/2021 - 11/2025}
{
    \item \textbf{Promedio académico: 4.6 / 5.0} (Mérito académico).
    \item Formación sólida en procesamiento digital de señales (DSP), modelado computacional, analítica numérica, automatización y comunicaciones.
    \item Desarrollo de proyectos enfocados en simulación de procesos, programación de algoritmos, validación experimental y reportes de ingeniería.
}

\newpage

\section{Proyectos Relevantes}

\cventry{Sistemas RAG y Agentes Autónomos}{AI Engineering \& Python}{Bogotá}{2026}
{
    \item Diseño e implementación de arquitecturas RAG (Retrieval-Augmented Generation) para consulta e indexación semántica sobre bases documentales técnicas y registros operativos.
    \item Desarrollo de flujos de agentes autónomos basados en LLMs con capacidad de ejecución de herramientas (tool calling), integración de APIs y orquestación multi-turno con LangChain.
    \item Implementación de protocolos de evaluación, guardrails y trazabilidad para monitorear la precisión de respuestas y asegurar la gobernanza de modelos en producción.
}

\cventry{DPS Analyzer \& Metrology AI (Series Temporales y MLOps)}{AM4}{Bogotá}{2026}
{
    \item Suite de instrumentación y diagnóstico analítico predictivo para infraestructura crítica: adquisición continua a 50 kHz, análisis espectral FFT y cálculo de índices de salud bajo normas IEEE/IEC.
    \item Despliegue contenerizado mediante Docker y automatización de pipelines CI/CD con GitHub Actions, optimizando el ciclo de entrega y monitoreo continuo de variables.
    \item Desarrollo de componentes backend para almacenamiento, consulta SQL, visualización y generación automática de reportes técnicos de alta precisión.
}

\section{Habilidades}
\begin{tcolorbox}[colback=lightgray, colframe=lightgray, boxrule=0pt, arc=6pt, left=10pt, right=10pt, top=10pt, bottom=10pt]
\noindent
\linespread{1.8}\selectfont
\skillbadge{Python (Avanzado)}
\skillbadge{IA Generativa \& LLMs}
\skillbadge{Sistemas RAG}
\skillbadge{Agentes Autónomos}
\skillbadge{LangChain / LlamaIndex}
\skillbadge{Databricks}
\skillbadge{MLOps \& CI/CD}
\skillbadge{Docker \& Contenedores}
\skillbadge{FastAPI \& REST APIs}
\skillbadge{SQL \& Modelado de Datos}
\skillbadge{Visión Computacional}
\skillbadge{PyTorch / Scikit-Learn}
\skillbadge{Linux SysAdmin}
\skillbadge{Pruebas QA (Pytest)}
\skillbadge{Gestión de Proyectos}
\end{tcolorbox}

\section{Idiomas}
\begin{itemize}[label={}, leftmargin=0.3cm, itemsep=6pt]
    \item \textbf{Español} \hfill Nativo \quad \textcolor{accent}{\faCircle \faCircle \faCircle \faCircle \faCircle}
    \item \textbf{Inglés} \hfill B2 (Profesional de Trabajo / Técnico) \quad \textcolor{accent}{\faCircle \faCircle \faCircle \faCircle} \textcolor{bulletline}{\faCircle}
    \item \textbf{Francés} \hfill A2 \quad \textcolor{accent}{\faCircle \faCircle} \textcolor{bulletline}{\faCircle \faCircle \faCircle}
\end{itemize}

\section{Certificaciones}
\begin{itemize}[label={\color{accentblue}\faAward}, leftmargin=0.6cm, itemsep=6pt]
    \item \textbf{Ethical Hacking \& Cybersecurity} --- Cisco Networking Academy \hfill \textit{04/2026}
    \item \textbf{CCNA: Introduction to Networks} --- Cisco Networking Academy \hfill \textit{07/2025}
    \item \textbf{Análisis Computacional de Datos Ambientales} --- Universidad Nacional de Colombia \hfill \textit{09/2025}
    \item \textbf{Introducción a la Fotónica Integrada y las Comunicaciones Cuánticas} --- UNAL \hfill \textit{07/2024}
    \item \textbf{Python Intermedio} --- SoloLearn \hfill \textit{03/2021}
\end{itemize}

\end{document}
